\section{Gluon self-energy diagrams}

\begin{figure}[b]
\centerline{\includegraphics[width=1.5in]{images/gluonselfl} \hspace{-5mm} \includegraphics[width=1.5in]{images/gluontriple}}
\caption{Gluon self-energy-type insertions into the right leg.
\label{gluself}}
\end{figure}

One may expect that
the plus-prescription form would appear after the addition of the
gluon self-energy diagrams, one of which is shown in Fig. \ref{gluself}a.
These diagrams have both the UV and collinear
divergences. On the lattice, the UV divergence is regularized by the lattice spacing.
In a continuum theory, one may use the
Polyakov prescription $1/z^2\to 1/(z^2 -a^2)$ for the gluon propagator.
The collinear divergences may be regularized by taking
a finite gluon mass $\lambda$. The result is a
$\ln (a^2 \lambda^2)$ contribution.
However, it does not have the $z$-dependence,
and apparently cannot help one to build the plus-prescription form for the
$\ln z_3^2$ part of the box contribution.

A possible way out is to represent $\ln (a^2 \lambda^2)$
as the difference $\ln (z_3^2 \lambda^2) - \ln (z_3^2/a^2 )$
of the evolution-type logarithm $\ln (z_3^2 \lambda^2)$
and a UV-type logarithm $\ln (z_3^2/a^2 )$.
The latter can be added
to the UV divergences of the diagrams \ref{linkself} and \ref{link}
corresponding to link self-energy and vertex corrections.
The $\ln (z_3^2 \lambda^2)$
part is then added to the evolution kernel.

To be on safe side with gauge invariance, we use
the dimensional regularization.
Then the analog of the $\ln (a^2 \lambda^2)$ logarithm
is a pole $1/(2-d/2)$ sometimes written as \mbox{$1/\epsilon_{\rm UV} - 1/\epsilon_{\rm IR}$. }
For our purposes, it is more convenient to
symbolically write it in a form similar to $\ln (a^2 \lambda^2)$.
Changing $\lambda \to\mu_{\rm IR}$ and $a \to 1/\mu_{\rm UV}$ we get $\ln (\mu_{\rm IR}^2/\mu_{\rm UV}^2)$,
and then split this into the difference $\ln (z_3^2 \mu_{\rm IR}^2) - \ln (z_3^2 \mu_{\rm UV}^2 )$.

We should also take into account the diagrams (one of them is shown in Fig. \ref{gluself}b)
with an extra gluon
line going out of the link-gluon vertex. The combined contribution
of the Fig. \ref{gluself} diagrams and their left-leg analogs is given by
\begin{align}
{g^2N_c\over 8\pi^2}
\frac1{2-d/2}
\left [2 - \frac{\beta_0}{2N_c} \right ]G_{\mu\alpha}(z)G_{\lambda \beta }(0) \ ,
\label{counter3}
\end{align}
where $\beta_0 =11N_c/3$ in gluodynamics, so that the terms in the square bracket combine into 1/6.
As discussed above, we will treat \mbox{$1/({2-d/2})$} as
\mbox{$\ln (z_3^2 \mu_{\rm IR}^2) - \ln (z_3^2 \mu_{\rm UV}^2 )$.}
